\documentclass[11pt,a4paper]{article}

\usepackage[spanish,es-noquoting]{babel}
\usepackage[utf8]{inputenc}
\usepackage[T1]{fontenc}
\usepackage{lmodern}
\usepackage{geometry}
\usepackage{booktabs}
\usepackage{tabularx}
\usepackage{array}
\usepackage{float}
\usepackage{enumitem}
\usepackage{graphicx}
\usepackage{xcolor}
\usepackage{hyperref}
\usepackage{fancyhdr}

\geometry{margin=2.5cm}
\setlength{\parskip}{0.6em}
\setlength{\parindent}{0pt}
\hypersetup{colorlinks=true, linkcolor=black, urlcolor=blue}

\setlength{\headheight}{15pt}

\pagestyle{fancy}
\fancyhf{}
\chead{\textbf{Universidad del Valle de Guatemala}}
\rhead{Semestre I 2026}
\cfoot{\thepage}
\renewcommand{\headrulewidth}{0.4pt}

\title{\textbf{Proyecto \#2 -- Data Mining (PBL)}}
\author{}
\date{}

\begin{document}

\maketitle

\begin{center}
\Large \textbf{Semestre I del 2026}
\end{center}

\hrule
\vspace{1cm}

\section*{🎯 Objetivo general}

Desarrollar un proyecto completo de minería de datos donde el estudiante:

\begin{itemize}[leftmargin=1.5cm]
\item Analice un conjunto de datos real
\item Formule un problema (sin que se le diga cuál)
\item Construya, compare y evalúe modelos
\item Justifique decisiones técnicas
\item Comunique resultados de forma profesional
\end{itemize}

\section*{🗂️ Conjunto de datos}

Se les entrega el archivo: \texttt{Datos.csv}

\hrule
\vspace{0.8cm}

\section*{📅 Semana 1 -- Comprensión del problema + EDA}

\subsection*{🎯 Objetivo}

Que los estudiantes descubran:

\begin{itemize}[leftmargin=1.5cm]
\item Qué representa el conjunto de datos
\item Qué problema pueden resolver
\item Primer análisis exploratorio
\end{itemize}

\subsection*{📌 Entregables}

\textbf{1. Documento breve (3--5 páginas):}

\begin{itemize}[leftmargin=1.5cm]
\item Descripción del conjunto de datos (variables, tipos)
\item Identificación del problema:
\begin{itemize}[leftmargin=0.8cm]
\item ¿Clasificación o regresión?
\item ¿Variable objetivo?
\end{itemize}
\item Análisis exploratorio:
\begin{itemize}[leftmargin=0.8cm]
\item Distribuciones
\item Correlaciones
\item Datos atípicos (Outliers)
\item Hipótesis iniciales
\end{itemize}
\end{itemize}

\textbf{2. Notebook (EDA):}

\begin{itemize}[leftmargin=1.5cm]
\item Visualizaciones (mínimo 5)
\item Comentarios interpretativos
\end{itemize}

\subsection*{📊 Rúbrica -- Semana 1 (20 pts)}

\begin{table}[H]
\centering
\begin{tabular}{lc}
\toprule
\textbf{Criterio} & \textbf{Puntos} \\
\midrule
Identificación correcta del problema & 5 \\
Calidad del EDA (visualizaciones + interpretación) & 7 \\
Comprensión de variables & 4 \\
Claridad del documento & 4 \\
\bottomrule
\end{tabular}
\end{table}

\hrule
\vspace{0.8cm}

\section*{📅 Semana 2 -- Preparación de datos + Modelos base}

\subsection*{🎯 Objetivo}

Construir primeros modelos sin optimización

\subsection*{📌 Contenido esperado}

\begin{itemize}[leftmargin=1.5cm]
\item Limpieza de datos (si aplica)
\item Selección de variables
\item Train/Test split
\item Modelos base (mínimo 3)
\end{itemize}

\subsection*{📌 Entregables}

\textbf{1. Notebook:}

\begin{itemize}[leftmargin=1.5cm]
\item Preprocesamiento
\item Modelos base
\item Métricas:
\begin{itemize}[leftmargin=0.8cm]
\item Exactitud (Accuracy)
\item Precisión, Recall, F1
\item Matriz de confusión
\end{itemize}
\end{itemize}

\textbf{2. Documento (2--3 páginas):}

\begin{itemize}[leftmargin=1.5cm]
\item Justificación de modelos elegidos
\item Comparación inicial
\end{itemize}

\subsection*{📊 Rúbrica -- Semana 2 (25 pts)}

\begin{table}[H]
\centering
\begin{tabular}{lc}
\toprule
\textbf{Criterio} & \textbf{Puntos} \\
\midrule
Correcto preprocesamiento & 6 \\
Implementación de modelos & 8 \\
Uso de métricas adecuadas & 6 \\
Interpretación de resultados & 5 \\
\bottomrule
\end{tabular}
\end{table}

\hrule
\vspace{0.8cm}

\section*{📅 Semana 3 -- Mejora de modelos + Selección}

\subsection*{🎯 Objetivo}

Optimizar modelos y seleccionar el mejor

\subsection*{📌 Contenido esperado}

\begin{itemize}[leftmargin=1.5cm]
\item Cross-validation
\item Afinación de Hiperparámetros (Hyperparameter tuning)
\item Comparación formal de modelos
\item Sobre ajuste vs Sub ajuste (Overfitting vs underfitting)
\end{itemize}

\subsection*{📌 Entregables}

\textbf{1. Notebook:}

\begin{itemize}[leftmargin=1.5cm]
\item Modelos optimizados
\item Comparación clara
\end{itemize}

\textbf{2. Documento (3--4 páginas):}

\begin{itemize}[leftmargin=1.5cm]
\item Justificación del modelo final
\item Discusión técnica:
\begin{itemize}[leftmargin=0.8cm]
\item Sesgo/varianza (Bias/variance)
\item Trade-offs
\end{itemize}
\end{itemize}

\subsection*{📊 Rúbrica -- Semana 3 (30 pts)}

\begin{table}[H]
\centering
\begin{tabular}{lc}
\toprule
\textbf{Criterio} & \textbf{Puntos} \\
\midrule
Uso correcto de validación & 8 \\
Optimización de modelos & 8 \\
Comparación rigurosa & 7 \\
Justificación técnica & 7 \\
\bottomrule
\end{tabular}
\end{table}

\hrule
\vspace{0.8cm}

\section*{📅 Semana 4 -- Interpretación + Presentación final}

\subsection*{🎯 Objetivo}

Comunicar resultados como un científico de datos

\subsection*{📌 Contenido esperado}

\begin{itemize}[leftmargin=1.5cm]
\item Interpretabilidad:
\begin{itemize}[leftmargin=0.8cm]
\item Importancia de características (Feature importance) (RF / árboles)
\item Coeficientes (logística)
\end{itemize}
\item Análisis de errores
\item Conclusiones
\item Recomendaciones
\end{itemize}

\subsection*{📌 Entregables}

\textbf{1. Informe final (8--12 páginas)}

Debe incluir:

\begin{itemize}[leftmargin=1.5cm]
\item Introducción
\item Metodología
\item Resultados
\item Discusión
\item Conclusiones
\end{itemize}

\textbf{2. Presentación (10--15 min)}

\begin{itemize}[leftmargin=1.5cm]
\item Problema
\item Enfoque
\item Modelo final
\item Resultados clave
\end{itemize}

\subsection*{📊 Rúbrica -- Semana 4 (25 pts)}

\begin{table}[H]
\centering
\begin{tabular}{lc}
\toprule
\textbf{Criterio} & \textbf{Puntos} \\
\midrule
Interpretación del modelo & 7 \\
Calidad del informe & 8 \\
Claridad de la presentación & 5 \\
Rigor técnico & 5 \\
\bottomrule
\end{tabular}
\end{table}

\hrule
\vspace{1cm}

\section*{🧮 Evaluación total}

\begin{table}[H]
\centering
\begin{tabular}{lc}
\toprule
\textbf{Semana} & \textbf{Puntos} \\
\midrule
Semana 1 & 20 \\
Semana 2 & 25 \\
Semana 3 & 30 \\
Semana 4 & 25 \\
\midrule
\textbf{Total} & \textbf{100} \\
\bottomrule
\end{tabular}
\end{table}

\section*{🧩 Reglas del proyecto}

\begin{itemize}[leftmargin=1.5cm]
\item Trabajo en grupos de 4 estudiantes
\item Uso libre de librerías (sklearn, pandas, etc.)
\item \textbf{Prohibido copiar notebooks completos}
\item Se evaluará \textbf{explicación, no solo código}
\end{itemize}

\vfill

\begin{center}
\textit{Semestre I del 2026}
\end{center}

\end{document}
